% =====================================================================
\section{The peaks, one by one}
\label{si:peakbypeak}
% =====================================================================

This section gives the full physical account of every feature in the
final atlas (\cref{fig:si-atlas}), channel by channel.  All values are
blind-census amplitudes at the consolidated operating point of
\cref{si:operating}, normalised within each channel, with the
$\omega_\tau=0$ pump--probe row removed.  Throughout, ``delay axis''
means the coherence that survives between the pulses (it never
radiates, so it is constrained only by what can be \emph{excited} and
what can \emph{store} phase), while ``detection axis'' means the
coherence that radiates into the analysed polarisation (constrained by
the emitting multipole).  This asymmetry does most of the explanatory
work below.

\subsection{What each channel can see}
\label{si:whatsees}

Detection is resolved by the \emph{magnetic} polarisation of the
emitted THz field, and for a given input geometry the analysed output
selects one Cartesian component of the total magnetic moment:
\begin{center}\small
\renewcommand{\arraystretch}{1.3}
\begin{tabular}{llll}
\toprule
geometry / analyser & detected moment & Fe content & Tm content\\
\midrule
A, cross ($H\!\parallel\!a$ in, $H\!\parallel\!c$ out) & $m_z$ &
  $F_z\equiv0$ (machine zero) & $\mu\lambda^2+\beta\lambda^1\lambda^2$\\
A, same ($H\!\parallel\!a$ out) & $m_x$ &
  $F_x$: carries $\qafm$ & $2.39\lambda^5+0.91\lambda^7$ ($\to\lambda^{4,6}$)\\
B, cross ($H\!\parallel\!c$ in, $H\!\parallel\!a$ out) & $m_x$ &
  $F_x$: carries $\qafm$ & $\lambda^{4\text{--}7}\equiv0$ (subalgebra)\\
B, same ($H\!\parallel\!c$ out) & $m_z$ &
  $F_z$: carries $\qfm$ & $\lambda^{4\text{--}7}\equiv0$\\
\bottomrule
\end{tabular}
\end{center}
Two entries in this table are exact statements rather than
approximations, and both do real work.  First, $F_z$ vanishes to
machine precision ($10^{-15}$ at every frequency) in geometry A: the
$\qafm$ magnon lives in $\delta F_x$ and has \emph{no} $c$-axis moment,
so the Fe sublattice is radiatively silent in the geometry-A cross
channel.  Anything seen there at a magnon frequency must therefore come
from the Tm ion.  Second, in geometry B the $E\!\parallel\!a$,
$H\!\parallel\!c$ drive closes on the $\{\lambda^1,\lambda^2,\lambda^3\}$
subalgebra, so the CEF coordinates $\lambda^{4\text{--}7}$ are
identically zero and both geometry-B channels are pure magnon maps.

\subsection{Geometry A, cross-polarised}
\label{si:peaksAcross}

\paragraph{$(\qafm,E_{12})=(0.90,\,0.50)$, amplitude $1.00$ --- the
magnetoelectric cross peak; measures $W^{yy}_1$.}
The $H\!\parallel\!a$ Zeeman field launches the $\qafm$ magnon, which
precesses through the delay at $0.90$\,THz.  It never radiates: it has
no $m_z$, and the analysed channel is $m_z$.  What it does instead is
modulate the Tm crystal field through the two-magnon vertex
$W^{yy}_1S_y^2\lambda^1$.  Because the vertex is quadratic in the
fluctuation and the two pulses each supply one magnon leg, the
$M_{\rm NL}$ subtraction keeps only the cross term, which
\emph{rectifies},
\begin{equation}
  h^1_{\rm NL}(t,\tau)=W^{yy}_1A^2\big[\cos(\qafm\tau)
  -\cos(\qafm(2t+\tau))\big]\theta(t),
\end{equation}
so the difference-frequency piece is a step force switched on at the
probe, carrying the delay phase $\cos(\qafm\tau)$ but no $t$
dependence.  Fed into the $E_{12}$ oscillator it produces
$M_{\rm NL}\propto\cos(\qafm\tau)[1-\cos\omega_{12}t]$: a factorised,
rank-one peak whose delay axis is the magnon and whose detection axis
is the CEF line.  The peak is thus magnetoelectric in the strict
sense---an E1-dark magnon made visible by borrowing the Tm dipole---and
its existence is a direct measurement of $W^{yy}_1$.
\emph{Energy bookkeeping:} $\delta S_y^2$ contains only DC and
$2\qafm$, so between free modes this vertex could never emit at
$\omega_{12}$; the $E_{12}$ quantum is drawn from the \emph{pulse edge}
(an intra-pulse difference-frequency, i.e.\ displacive, process).  This
is verified rather than asserted: stretching the pulse from
$\sigma=0.25$ to $1.0$ at fixed amplitude collapses the peak by $4.5$
decades ($51.7\to1.7\times10^{-3}$), the adiabatic suppression a
displacive mechanism demands, while switching the direct Tm drive off
changes it by $<10^{-4}$.
\emph{Scaling:} $\propto A_{\rm Fe}^2$ (two magnon legs), independent
of the Tm drive.  \emph{Temperature:} flat below $\sim10$\,K, since
both the excitation and the storage run through the long-lived magnon.
\emph{Linewidth along $\omega_t$:} the CEF value, $0.038$\,THz.

\paragraph{$(E_{12},2E_{12})=(0.50,\,1.00)$, amplitude $0.97$ ---
second-harmonic emission; measures the Tm hyperpolarisability $\beta$.}
The pump stores an $E_{12}$ coherence, which sets the delay axis at
$0.50$\,THz.  The detection axis is its \emph{second harmonic},
$2E_{12}=1.0000$\,THz.  A moment linear in the qutrit coordinates
cannot radiate a harmonic, but the Tm$^{3+}$ ion occupies the $4c$
site, whose point symmetry is a single mirror with \emph{no inversion}
(Fe, at $4b$, sits on an inversion centre and is therefore excluded
from this mechanism).  The lowest allowed addition to the emitted
moment is
\begin{equation}
  m_z=\mu\,\lambda^2+\beta\,\lambda^1\lambda^2 ,
\end{equation}
which is $\mathcal T$-odd as a magnetisation must be ($\lambda^1$
$\mathcal T$-even, $\lambda^2$ $\mathcal T$-odd) and, since both factors
oscillate at $\omega_{12}$, radiates at $2\omega_{12}$.  It is a
detection-side (hyperpolarisability) term: the dynamics is untouched.
Evaluated on the trajectories, this channel is strikingly clean---it
contributes essentially one peak, $5.8\times$ larger than its own
$(E_{12},E_{12})$ diagonal, with only a weak $(\qafm,2E_{12})$
companion, and it is identical in the local and global sublattice
frames.  Calibrating $\beta$ by the observed intensity parity with the
main peak gives $\beta\simeq7\mu$.
\emph{Discriminators} (this is the peak whose assignment matters most):
the line sits at exactly $2E_{12}$ and \emph{tracks $E_{12}$}, so under
temperature or field it follows the CEF level and does \emph{not}
soften toward the spin reorientation as any magnon-derived feature
would; it is one field order above $(\qafm,E_{12})$ and therefore grows
faster with fluence; its width is twice the $E_{12}$ linewidth,
$\approx0.076$\,THz, against $0.003$\,THz for anything magnon-derived;
and it is independent of both $W^{xz}_1$ and the Zeeman drive.

\paragraph{$(-\qafm,E_{12})$, amplitude $0.52$ --- the rephasing
partner of the main peak.}
Not a separate mechanism.  The $W$ vertex is symmetric under the two
time orderings of the magnon coherence, so the signal appears on both
signed sides of the delay axis, with the non-rephasing branch stronger.
Experiments that display folded $|\omega_\tau|$ merge this with the
main peak; experiments that resolve the sign should see it at roughly
half amplitude.

\paragraph{$(\qafm,2\qafm)=(0.90,\,1.79)$, amplitude $0.39$ --- a
standing prediction.}
The same hyperpolarisability term that produces the $2E_{12}$ peak must
also double any \emph{other} frequency present in the Tm dipole.  The
magnon-forced component (next entry) oscillates at $\qafm$, so
$\beta\lambda^1\lambda^2$ radiates its harmonic at $2\qafm=1.79$\,THz
on the $\qafm$ delay row.  Observing it would confirm the
hyperpolarisability assignment independently of the $2E_{12}$ peak;
its absence, with $(E_{12},2E_{12})$ present, would falsify
Eq.~(\ref{eq:shg}) in favour of some other harmonic mechanism.  It may
lie outside the detection bandwidth.

\paragraph{$(\qafm,\qafm)=(0.90,\,0.89)$, amplitude $0.36$ --- the
magnon diagonal, seen through the Tm ion.}
This one is subtle, because $F_z\equiv0$ means the magnon cannot
radiate here at all.  What radiates is the \emph{Tm} $m_z$ dipole,
which carries a small component oscillating at the magnon frequency.
The mechanism is an off-resonant forced response: in $\Gamma_2$ the
static vertex $W^{xz}_1S_xS_z\lambda^1$ contains
$G_z\,\delta S_x\,\lambda^1$, a term \emph{linear} in the magnon
fluctuation, which shakes the Tm oscillator at $0.90$\,THz.  A driven
oscillator responds at the driving frequency, not its own, so the
$0.50$\,THz Tm dipole acquires a $0.90$\,THz wobble
(\cref{fig:si-forced}).  This is verified in detail: the sideband sits
at $0.8999$\,THz (the magnon, not the $\qfm+E_{12}=0.876$
combination); its amplitude is \emph{linear} in the Zeeman drive
($\times1.70$ for a $\times1.67$ increase, tracking the source at
$\times1.65$) whereas a harmonic would be quadratic; the transfer
coefficient $|\lambda^2|_{0.90}/|m_x|_{0.90}=0.146$ is
drive-independent, the hallmark of a linear susceptibility
$\chi\propto\omega_{12}/(\omega_{12}^2-\omega_q^2)$ evaluated off
resonance; and the admixture grows linearly with $W^{xz}_1$
($0.09\%\to1.2\%\to2.6\%$ at $W^{xz}_1=0/0.01/0.02$).
The same mechanism, with an $E_{12}$ delay instead of a $\qafm$ delay,
gives the weak $(E_{12},\qafm)$ reverse-transfer feature; being
doubly-$W^{xz}_1$-mediated (transfer $\times$ detection) its amplitude
is \emph{quadratic} in $W^{xz}_1$ ($3.7\times$ measured against
$4\times$ expected), making it the atlas's dedicated $W^{xz}_1$ meter.
Because the geometry-B mixer sweep bounds $W^{xz}_1\lesssim0.02$, the
model caps this feature well below the main peak---which is why the
reported second peak is assigned to the harmonic instead.

\paragraph{$(\pm E_{12},E_{12})$, amplitudes $0.33$ and $0.32$ --- the
one-quantum diagonal.}
The directly driven $E_{12}$ coherence, read out through the same
$m_z$ dipole.  Diagonals are present in every 2D map and are not
cross peaks; they are reported here for completeness.  Their size
relative to the cross peaks is a sensitive fluence diagnostic: the
diagonal falls from $1.45$ of the main peak at $A_{\rm Fe}=1.2$ to
$0.16$ at $0.12$ and $0.06$ at $0.04$, which is how the operating point
was fixed (\cref{si:operating}).

\subsection{Geometry B, cross-polarised}
\label{si:peaksBcross}

Here the analysed moment is $m_x$, which the $\qafm$ magnon carries
directly, so both peaks emit at $0.90$\,THz through the Fe magnetic
dipole and the Tm ion is radiatively absent.

\paragraph{$(\qfm,\qafm)=(0.38,\,0.90)$, amplitude $1.00$ --- the
magnon--magnon peak; measures the Fe sector alone.}
$H\!\parallel\!c$ launches the quasi-ferromagnetic magnon, which
precesses at $0.38$\,THz through the delay; single-ion anisotropy and
the Dzyaloshinskii--Moriya interaction couple the two magnon branches,
so the stored $\qfm$ phase is re-emitted on the $\qafm$ branch.  No
Fe--Tm coupling is involved at any vertex, which is why this peak is
temperature-flat: it is a pure two-magnon process, independent of the
qutrit populations.

\paragraph{$(E_{12},\qafm)=(0.53,\,0.89)$, amplitude $0.45$ --- the
gated transfer peak; measures $\kappa^B_{1y}$.}
$E\!\parallel\!a$ stores a Tm $E_{12}$ coherence through the delay; a
Fe--Tm vertex then transfers it to the magnon, which radiates via M1.
Isolating the vertices settles which one: with $\kappa^B_{1y}$ alone
the peak is present, with the static hybridisation alone it is absent,
and \emph{no strength of the static vertex rescues it}.  A sweep over
$W^{xz}_1\in\{0.01,\dots,0.09\}$ shows the static vertex never creates
a second resolved peak; instead the two delay rows coalesce and the
leading peak's delay phase migrates continuously
($0.38\to0.40\to0.48$\,THz), while hybridisation-shifted $E_{12}$
emission leaks in at $0.44$--$0.45$\,THz.  This is the structural
signature of a mixer: a static vertex \emph{moves} peaks (level
repulsion, which also shifts the modes off their 1D-measured
positions), whereas a gated converter \emph{adds} a peak at unshifted
lines.  The experiment shows two resolved peaks at the unshifted
energies, so the gate is required.  Physically $\kappa^B$ supplies the
$0.50\to0.90$\,THz mismatch impulsively from the photon while
$B_z(t)\neq0$, and the very same gate makes it \emph{identically
silent} in geometry A ($B_z=0$), where its background would otherwise
contaminate the clean $W^{yy}_1$ peak.  \emph{Temperature:} falls as
the population difference $n_1-n_2$
($0.46\to0.45\to0.40\to0.36$ at $0/5/8/10$\,K), because it runs off the
$E_{12}$ coherence amplitude---the sharpest thermal discriminator in
the atlas against the temperature-flat magnon--magnon peak.

\paragraph{$(-\qfm,\qafm)$ at $0.17$ and $(\qafm,\qafm)$ at $0.14$.}
The rephasing partner of the magnon--magnon peak, and the $\qafm$
diagonal.  Note the diagonal here is genuine Fe M1 emission, unlike its
geometry-A counterpart.

\paragraph{$(E_{12}{+}\qafm,\qafm)=(1.43,\,0.90)$, amplitude $0.07$.}
A combination tone: the sum-frequency of the two stored coherences on
the delay axis, radiating on the magnon line.  It is $W$-fed and
therefore scales with the same coupling as the transfer peak; its
observation would provide an independent handle on the $W$ scale.

\paragraph{The post-drive buildup.}
The measured geometry-B signal continues to grow after the pulse has
passed.  No coherent vertex can do this once the measured linewidths
are enforced: because $B\!\parallel\!c$ never drives $\qafm$ directly,
all geometry-B magnon amplitude is \emph{converted}, and a converted
signal inherits its source's lifetime---with $T_2(E_{12})=8$\,ps
everything is over within $\sim10$\,ps, and the artefact-free Gabor
analysis shows the amplitude peaking at pulse end in every vertex
configuration.  Even the $\tau$-labelled population grating fails: a
full 2DCS run with the population channel on leaves the
$\omega_\tau=E_{12}$ row identical to the baseline, because the
grating's delay label is written by the $E_{12}$ coherence (hence
$T_2$-confined) and its force is an impulsive step.  The buildup
therefore requires the one reservoir that outlives every $T_2$: the CEF
\emph{populations}, fed incoherently by the dissipated pulse energy and
acting on the magnon through the population channel
$W_3\,S_\alpha S_\beta\lambda^3$---the mechanism of thermally driven
spin reorientation in this material.  Implemented in the solver
(dissipated coherence energy shifts the $\lambda^3$ equilibrium; the
$\Gamma_3$ relaxation chases it, so the rise time is $1/\Gamma_3$ by
construction), it reproduces the observation: the $\omega_\tau=E_{12}$
row falls off the impulsive peak, recovers on the repopulation time,
and persists beyond $100$\,ps, $50$--$100\times$ above the
coherent-only model at late times.  The recovery time measures the Tm
repopulation time and the late-to-impulsive ratio calibrates
$c_{\rm heat}W_3$.

\subsection{Geometry A, same-polarised}
\label{si:peaksAsame}

The analysed moment is $m_x$, which receives the Tm $a$-axis moments
($\lambda^{5,7}$, equivalently their quadratures $\lambda^{4,6}$) and
the Fe $F_x$.  The organising fact of this channel is a selection rule.

\paragraph{The $z$-dark $1\to3$ dipole.}
The measured inventory contains delay frequencies
$\omega_\tau\in\{E_{12},\qafm\}$ only---there is no feature anywhere
with delay phase $E_{13}$.  A bright $1\!\to\!3$ dipole would
necessarily store an $E_{13}$ coherence at first order and produce an
$(E_{13},E_{23})$ peak with \emph{exactly} the same dipole product
$\mu_{12}\mu_{13}\mu_{23}$ as the observed $(E_{12},E_{23})$ peak, so
no parameter choice could separate them.  The absence of the whole row
is therefore a selection rule, $\mu^z_{13}\simeq0$: the natural
consequence of a parity-alternating CEF ladder ($\ket1,\ket3$ even,
$\ket2$ odd under the site mirror), for which the $z$-dipole connects
$1\!\leftrightarrow\!2$ and $2\!\leftrightarrow\!3$ but not
$1\!\leftrightarrow\!3$.  This does not contradict the measured
$E_{13}$ oscillator strength ($f=8.7$), which lives in the orthogonal
polarisations: the same transition is bright in one polarisation and
dark in another, a fact the 2D map resolves and 1D absorption---summed
over its bright polarisations---does not.  The residual admixture is
not assumed but \emph{measured}: it multiplies the $E_{13}$ emission
only, so the observed ordering of the two transfer peaks fixes it at
$\mu^z_{13}/\mu^z_{23}=0.14\%$.

\paragraph{$(E_{12},E_{23})=(0.49,\,0.70)$, amplitude $1.00$ --- the
dominant transfer, via a population.}
With the single-action route removed by the dark dipole, the dominant
pathway acts twice: one probe interaction converts the stored $E_{12}$
coherence into a level-2 \emph{population} that still carries the
$\omega_\tau=E_{12}$ phase, and a second converts that population into
the emitting $E_{23}$ coherence.  Being co-rotating, this route lands
\emph{non-rephasing}---the observed side---and it involves $\mu_{13}$
at no vertex.  It is one field order higher than the $E_{13}$ partner,
which is why it dominates only above a threshold fluence and why the
hierarchy inverts at weak drive; that sensitivity is what fixes the
operating point jointly with the cross channel.

\paragraph{$(E_{12},E_{23})$ at zero detection frequency --- the
rectified line.}
Taken at its zero-frequency output, the same double action rectifies
the stored coherence into a slow net magnetisation of the detection
window.  $S_{\rm DC}(\tau)=\langle M_{\rm NL}\rangle_t$, after cubic
drift removal, shows a sharp line at exactly $+E_{12}$---the
$(E_{12},0)$ entry of the measured list, and the sharpest feature of
the channel.  It is the same storage physics with a DC output leg.

\paragraph{$(E_{12},E_{13})=(0.49,\,1.20)$, amplitude $0.80$ --- the
throttled partner.}
A \emph{single} probe action on the bright $2\!\leftrightarrow\!3$
dipole rotates the bra to give $\rho_{13}$, emitting at
$E_{13}=1.20$\,THz; the route is again co-rotating, hence the same
non-rephasing side, as observed.  Although lower order than the
$E_{23}$ route, its emission passes through the nearly dark
$\mu^z_{13}$, which caps it below the dominant peak.  The measured
ratio of the two transfer peaks is thus a direct measurement of the
residual $1\!\to\!3$ admixture.

\paragraph{$(-E_{12},E_{23})$, amplitude $0.67$ --- the rephasing
remnant; a hard prediction.}
The mirror of the dominant transfer.  Because the single-ion dynamics
is quantum-exact, this was checked against a standalone qutrit
integrator: the rephasing partner retains $0.47$--$0.58$ of the
dominant peak across all drive strengths, dipole-phase choices and
dampings.  It is an exact floor for any three-level ion under this
pulse.  If the experimental map truly contains nothing at
$(-E_{12},E_{23})$, the suppression must come from outside the
point-ion model---quadrant folding in the display convention, or
propagation and phase-matching in a thick sample.

\paragraph{$(E_{13},E_{23})=(1.20,\,0.71)$, amplitude $0.32$ --- the
two-quantum row.}
An $E_{13}$ delay coherence does survive the dark dipole, but only at
\emph{second} order: the pump drives $1\!\to\!2$ and $2\!\to\!3$ in
succession, building $\rho_{13}$ through the ladder rather than
directly.  Being second order it is far weaker than the first-order
row a bright $\mu^z_{13}$ would have produced, which is why the model
is consistent with an inventory that lists no $E_{13}$ delay features
while still predicting a weak one.  The analogous $2E_{12}$
two-quantum line, absent from the exact single-ion dynamics, is a
genuine inter-site (mean-field) product.

\paragraph{$(\qafm,E_{13})=(0.90,\,1.20)$, amplitude $0.08$ ---
thermally activated.}
The magnon is excited by the Zeeman channel and sets the delay phase;
the Fe--Tm couplings hand the modulation to the Tm ion, which emits
weakly at $E_{13}$ through the throttled $\mu^z_{13}$.  It grows
steeply with the thermal population of level 2
($0.10\to0.17\to0.51$ relative at $0/10/20$\,K), the expected trend for
a route assisted by excited-level population, and is therefore the
channel's best thermometer.

\paragraph{The missing $(\qafm,\qfm)=(0.90,\,0.38)$.}
The one listed feature the model does not place in this channel.  The
carrier exists---$\qfm$ emission is a $b$/$c$-axis affair, as
$\Gamma_2$ requires---but the analysed moment here is $m_x$, and the
$\qfm$ mode radiates through $m_z$, which is dark in geometry A at the
$10^{-11}$ level.  The same transfer is the \emph{dominant} peak of the
geometry-B same-polarised channel (below), so a weak
$(\qafm,\qfm)$ in the geometry-A same-polarised map is naturally
polarisation leakage of that channel.  Discriminator: the feature
should scale as the square of the analyser leakage angle in a
polariser-purity scan.

\subsection{Geometry B, same-polarised: a pure magnon map (prediction)}
\label{si:peaksBsame}

No data exist for this channel yet, so the entire census is a
prediction, and a clean one: the analysed moment is $m_z$, and because
the geometry-B drive closes on the $\{\lambda^{1,2,3}\}$ subalgebra the
Tm CEF coordinates are \emph{identically} zero.  This is therefore a
pure Fe magnon map, and observing \emph{any} CEF emission in it would
falsify the subalgebra structure of the drive.
\begin{center}\small
\renewcommand{\arraystretch}{1.25}
\begin{tabular}{lll}
\toprule
peak & amplitude & content\\
\midrule
$(\qafm,\qfm)=(0.90,0.38)$ & $1.00$ & magnon--magnon transfer, reverse
  direction of the geometry-B cross peak\\
$(E_{12},\qfm)=(0.52,0.38)$ & $0.42$ & Tm$\to$magnon transfer read out
  on the $\qfm$ branch\\
$(\qfm,\qfm)$ & $0.41$ & $\qfm$ diagonal\\
$(-\qfm,E_{12})$ & $0.38$ & rephasing partner with $E_{12}$ emission\\
\bottomrule
\end{tabular}
\end{center}
The headline is that $(\qafm,\qfm)$ dominates here---precisely the
feature absent from the geometry-A same-polarised channel---which makes
the two channels a matched pair and turns the leakage hypothesis above
into a quantitative test.
